
\subsection*{7.16. Passage aux fibres géométriques.}
Soit \(\bar K\) une clôture algébrique de \(K=k(t)\), limite inductive de ses sous-extensions finies \(K_i\), et soit \(\bar A=\lim \bar A_i\) le normalisé de \(A\) dans \(\bar K\), limite inductive de ses normalisés \(\bar A_i\) dans les \(K_i\). Choisissons un idéal maximal \(\bar{\underline{m}}\) de \(\bar A\), induisant donc sur chaque \(\bar A_i\) un idéal maximal \(\underline{m}_i\), et soit \(B_i=A_{i\,\underline{m}_i}\), qui par Krull--Akizuki est encore un anneau de valuation discrète, dont nous désignons le spectre par \(S_i\), les points par \(s_i\) et \(t_i\). Le corps résiduel \(\bar k=\bar A/\bar{\underline{m}}\) est une clôture algébrique de \(k=k(s)\), et c'est la limite inductive des corps résiduels \(k_i=A_i/\underline{m}_i=k(s_i)\) des \(B_i\). Soit \(\bar t=\Spec(\bar K)\), \(\bar s=\Spec(\bar k)\). \emph{Supposons maintenant} \(X\) \emph{de type fini sur} \(S\), donc \(X_{\bar t}\) et \(X_{\bar s}\) de type fini sur \(\bar t\) resp. \(\bar s\), donc \emph{noethériens}. En vertu de IV 3.2.3, on a
\[
  K_\bullet(X_{\bar t})\xleftarrow{\sim}\varinjlim_i K_\bullet(X_{t_i}),
  \qquad
  K_\bullet(X_{\bar s})\xleftarrow{\sim}\varinjlim_i K_\bullet(X_{s_i}).
\]
D'autre part, en vertu de 7.15 on a un homomorphisme du système inductif des \(K_\bullet(X_{t_i})\) dans le système inductif des \(K_\bullet(X_{s_i})\). Passant à la limite, on trouve un homomorphisme canonique :
\begin{equation}\tag{7.16.1}
  K_\bullet(X_{\bar t})\longrightarrow K_\bullet(X_{\bar s}),
\end{equation}
appelé encore, comme de juste, \emph{homomorphisme de spécialisation}. Il induit un homomorphisme \(\Gr_\bullet(X_{\bar t})\longrightarrow \Gr_\bullet(X_{\bar s})\) grâce à 7.4, en utilisant le fait (dont la vérification est laissée au lecteur) que le foncteur \(\Gr_\bullet\) (tout comme \(K_\bullet\)) commute aux limites qui interviennent ici. De même, lorsque \(f\) est lisse, on peut interpréter (7.16.1) comme un homomorphisme sur les \(K^\bullet\), qui sera d'ailleurs un homomorphisme d'anneaux, et même un homomorphisme de \(\lambda\)-anneaux si \(X\) est séparé~; dans ce cas, il induit donc un homomorphisme sur les \(\Gr^\bullet\).

